%%% ---------- 导言区 ---------- %%%
\usepackage[a4paper,left=2.5cm,right=2.2cm,top=2.5cm,bottom=2.5cm]{geometry}

% 数学
\usepackage{amsmath,amssymb,amsfonts,bm}
% 西文与数学字体：Times 系（与原文 TimesNewRomanPSMT 一致）
% 优先使用 newtx（完整 TeX Live / Overleaf 均自带）：Times 正文 + Times 数学；
% 若无 newtx，则退回 TeX Gyre Termes（Times 的自由克隆）作正文；
% 若两者都没有，保持 LaTeX 默认字体，文档仍可正常编译。
\IfFileExists{newtxtext.sty}{%
  \usepackage{newtxtext}%
  \usepackage{newtxmath}%
}{%
  \IfFileExists{texgyretermes-regular.otf}{%
    \setmainfont{texgyretermes}[%
      Extension    = .otf,
      UprightFont  = *-regular,
      BoldFont     = *-bold,
      ItalicFont   = *-italic,
      BoldItalicFont = *-bolditalic]%
  }{}%
}

% 图表
\usepackage{graphicx}
\usepackage{booktabs}
\usepackage{array}
\usepackage{multirow}
\usepackage{caption}
\usepackage{float}

% 列表与杂项
\usepackage{enumitem}
\usepackage{xcolor}
\usepackage{listings}
\usepackage{url}
\usepackage{pifont}

\graphicspath{{figures/}}

% 图表标题：加粗标签 + 冒号，与原文一致
\captionsetup{labelsep=colon,labelfont=bf,font=small,skip=6pt}

% 章节格式：一、二、…（居中四号黑体）；小节用 5.1 / 5.1.1 的阿拉伯编号
\ctexset{
  section     = {name={,、}, number=\chinese{section},
                 format=\centering\bfseries\zihao{4}, beforeskip=18pt plus 3pt minus 3pt,
                 afterskip=12pt plus 2pt minus 2pt, aftername=},
  subsection  = {number=\arabic{section}.\arabic{subsection},
                 format=\raggedright\bfseries\zihao{-4}, aftername=\hspace{1em},
                 beforeskip=12pt plus 2pt minus 2pt, afterskip=8pt plus 1pt minus 1pt},
  subsubsection={number=\arabic{section}.\arabic{subsection}.\arabic{subsubsection},
                 format=\raggedright\bfseries\zihao{5}, aftername=\hspace{1em},
                 beforeskip=10pt plus 2pt minus 2pt, afterskip=6pt plus 1pt minus 1pt},
}

% 行距
\linespread{1.3}

% 常用记号
\newcommand{\disk}{\mathit{disk}}
\newcommand{\vect}[1]{\mathbf{#1}}
\newcommand{\degr}{\ensuremath{^{\circ}}}
\DeclareMathOperator{\cross}{cross}
\DeclareMathOperator{\slack}{slack}
\DeclareMathOperator{\diam}{diam}
\DeclareMathOperator*{\maxGap}{max\,Gap}
\DeclareMathOperator{\vertop}{vert}

% 代码清单样式（Python）
\definecolor{lstkw}{RGB}{0,0,160}
\definecolor{lststr}{RGB}{160,0,0}
\definecolor{lstcmt}{RGB}{0,120,0}
\definecolor{lstnum}{RGB}{110,110,110}
\lstdefinestyle{pystyle}{
  language=Python,
  basicstyle=\ttfamily\scriptsize,
  keywordstyle=\color{lstkw}\bfseries,
  stringstyle=\color{lststr},
  commentstyle=\color{lstcmt}\itshape,
  numbers=left,
  numberstyle=\rmfamily\tiny\color{lstnum},
  numbersep=8pt,
  stepnumber=1,
  breaklines=true,
  breakatwhitespace=false,
  breakindent=0pt,
  postbreak=\mbox{\textcolor{lstnum}{$\hookrightarrow$}\space},
  columns=fullflexible,
  keepspaces=true,
  showstringspaces=false,
  extendedchars=true,
  tabsize=4,
  xleftmargin=2.2em,
  framexleftmargin=2.2em,
  aboveskip=6pt,
  belowskip=6pt,
}
\lstset{style=pystyle}
